\documentclass[11pt]{article}
\usepackage[margin=1in]{geometry}

\usepackage{graphicx}
\usepackage{float}
\usepackage{booktabs}
\usepackage{tabularx}
\usepackage{ragged2e}

\title{Reporte de trabajo en equipo --- Entrega 2 (1 página)}
\author{}
\date{\today}

\begin{document}
\maketitle

% Evita cortes de palabras (especialmente nombres/rutas)
\hyphenpenalty=10000
\exhyphenpenalty=10000

\section*{Roles y contribuciones}

{\scriptsize
\setlength{\tabcolsep}{3pt}
\renewcommand{\arraystretch}{2.0}

\begin{table}[H]
\centering
\begin{tabularx}{\textwidth}{
>{\RaggedRight\arraybackslash}p{4.2cm}
>{\RaggedRight\arraybackslash}p{2.6cm}
>{\RaggedRight\arraybackslash}X
>{\RaggedRight\arraybackslash}p{3.2cm}}
\toprule
\textbf{Miembro} & \textbf{Rol} & \textbf{Entregables} & \textbf{Evidencia} \\
\midrule
David Alan Delgado Zarco & Modelado y metodología &
Label sin leakage, features, split temporal, modelos, tuning de umbral &
Commits notebook, outputs (tablas/figuras) \\
Octavio Esquivel Alvarez del Castillo & Reporte (LaTeX) + QA &
Plantilla LaTeX ($\leq$10 pág), integración de evidencias, revisión de rúbrica &
Commits report/, PDF final \\
John Williams Morrill & Tablero (Streamlit) &
Dashboard funcional, filtros, export, capturas del tablero &
dashboard/streamlit\_app.py, evidence/dashboard \\
Jason Rosher Morrill & MLOps / Infra (EC2+MLflow) &
MLflow en EC2 estable, persistencia, runs, capturas MLflow/EC2 &
infra/ + evidence/aws y evidence/mlflow \\
\bottomrule
\end{tabularx}
\end{table}
}

\end{document}